\section{Mitigating Scarcity and Imbalance}
\label{sec:scarcity}

Data scarcity in rare disease research manifests in multiple dimensions beyond simple sample counts. Small cohorts exhibit limited phenotypic variability, constraining model exposure to the full spectrum of disease presentations. Class imbalance, where rare disease cases constitute a small fraction relative to healthy controls or common conditions, further compounds training challenges. This section examines strategies through which generative models address these interconnected challenges, with emphasis on quantitative expansion rates, quality assessment methods, and approaches to preserving biological plausibility.

\subsection{Dataset Expansion Strategies}

The most direct application of generative models addresses absolute sample scarcity through synthetic data creation. Finetti et al. analyzed expansion rates across 118 studies, finding that medical imaging datasets achieved an average threefold increase through classical and deep generative augmentation \cite{finetti2025data}. However, this aggregate metric masks substantial variability across applications and techniques. Frid-Adar et al. expanded an initial dataset of 182 liver lesion CT scans through classical augmentation (rotations, translations, flips, scaling) to approximately 30,000 augmented samples per lesion class before applying GAN synthesis \cite{fridadar2018gan}. The combination of classical and deep generative approaches enabled training convolutional neural networks that would otherwise fail to converge on the original small dataset.

The choice of expansion magnitude depends on balancing dataset diversity against computational feasibility and the risk of introducing artifacts. Excessive reliance on synthetic data can lead to models that perform well on generated samples but fail to generalize to real-world variation. Mendes et al. documented synthetic cohort expansions ranging from threefold to tenfold across multiple rare disease applications \cite{mendes2025synthetic}. For acute myeloid leukemia studies, CTAB-GAN+ generated synthetic patient records that tripled the available training data while maintaining the statistical characteristics of rare genetic variants. The key distinction from classical augmentation lies in the generative models' capacity to produce novel combinations of features rather than simple transformations of existing samples.

Temporal and sequential data present unique expansion challenges due to the need for maintaining realistic longitudinal patterns. Zhong et al. addressed this through conditional generation in MedDiffusion, where the model generates synthetic patient visits based on medical history rather than independently \cite{zhong2024meddiffusion}. This approach ensures that synthetic disease trajectories exhibit plausible progression patterns consistent with clinical understanding. The model generated synthetic visits for patients with kidney disease, chronic obstructive pulmonary disease (COPD), and amnesia, expanding datasets with 2,810 to 7,314 positive cases by factors ranging from two to five depending on the specific prediction task requirements.

\subsection{Quality Assessment and Validation}

A critical challenge in synthetic data generation is establishing appropriate quality metrics beyond downstream model performance. Finetti et al. identified that over 40\% of reviewed studies provided insufficient quantitative detail regarding dataset sizes and expansion rates, while fewer than 33\% reported both pre- and post-augmentation metrics \cite{finetti2025data}. This lack of standardized reporting complicates comparative analysis and reproducibility assessment across studies.

Visual assessment by domain experts has emerged as a common validation approach for medical imaging applications. Frid-Adar et al. had radiologists evaluate synthetic liver lesions for anatomical plausibility and diagnostic utility \cite{fridadar2018gan}. The experts confirmed that GAN-generated lesions exhibited realistic texture patterns and morphological characteristics across all three lesion classes (cysts, metastases, hemangiomas). However, this validation method scales poorly and introduces subjectivity, as inter-rater reliability may vary and systematic evaluation protocols remain underdeveloped.

Statistical validation through distribution matching provides more objective assessment. Choi et al. demonstrated that medGAN-generated electronic health records preserved the marginal distributions of diagnosis and medication codes while maintaining appropriate correlations between co-occurring conditions \cite{choi2017generating}. The synthetic data reproduced the prevalence of rare conditions within the expected statistical variation of the real dataset. Tests comparing mean, variance, and higher-order moments across real and synthetic distributions confirmed statistical fidelity. Privacy assessment through membership inference attacks and attribute disclosure metrics further validated that synthetic records did not memorize individual patients from the training data.

Performance-based validation evaluates synthetic data utility through downstream prediction tasks. Frid-Adar et al. measured liver lesion classification performance with and without synthetic augmentation, observing improvements from 78.6\% to 85.7\% sensitivity and from 88.4\% to 92.4\% specificity \cite{fridadar2018gan}. These 7.1 and 4.0 percentage point gains demonstrate that synthetic data meaningfully enhanced model generalization. Similarly, Mendes et al. reported that combining synthetic and real brain MRI data achieved 85.9\% classification accuracy, with Dice score improvements ranging from 3\% to 15\% for segmentation robustness \cite{mendes2025synthetic}. Zhong et al. showed that MedDiffusion improved health risk prediction PR-AUC from 71.05\% to 77.88\% for kidney disease, representing a 6.83 percentage point gain over the best non-augmented baseline \cite{zhong2024meddiffusion}.

\subsection{Addressing Class Imbalance}

Class imbalance, where rare disease cases are substantially outnumbered by controls, presents distinct challenges from overall data scarcity. Traditional oversampling techniques such as SMOTE (Synthetic Minority Over-sampling Technique) generate synthetic minority class examples through interpolation between existing samples. While effective for tabular data with limited dimensionality, SMOTE struggles with high-dimensional medical imaging and complex sequential records. Deep generative models offer more sophisticated alternatives by learning the underlying distribution of the minority class and generating novel examples that extend beyond linear interpolation.

Conditional generation mechanisms enable targeted synthesis of underrepresented disease subtypes. Conditional GANs and VAEs accept class labels as additional inputs, directing the generator to produce samples from specific diagnostic categories. This proves particularly valuable in rare disease contexts where certain phenotypic variants or genetic subtypes may have fewer than ten documented cases. Finetti et al. documented applications in hematological malignancies, where acute lymphoblastic leukemia subtypes exhibit highly imbalanced class distributions \cite{finetti2025data}. Synthetic generation of underrepresented subtypes enabled training classifiers that would otherwise collapse to predicting only the majority class.

The effectiveness of class balancing through synthetic generation depends critically on the original training set containing sufficient minority class examples to capture distributional characteristics. When fewer than 20-30 examples exist, even sophisticated generative models risk producing samples that merely reproduce training set variation rather than capturing the true population distribution. In such extreme scarcity scenarios, transfer learning from related conditions or incorporating prior knowledge through rule-based constraints may provide necessary inductive biases. Mendes et al. described approaches combining synthetic generation with domain expert input to ensure generated samples respect known biological constraints \cite{mendes2025synthetic}.

\subsection{Biological Plausibility and Artifact Detection}

A persistent concern with synthetic biomedical data is the potential generation of biologically implausible samples that, while statistically consistent with training data, violate known physiological or anatomical principles. Generative models trained purely on data-driven objectives lack explicit mechanisms to enforce domain constraints. A GAN might generate a liver lesion with impossible tissue density values or an electronic health record with contradictory diagnosis codes that never co-occur in clinical practice.

Several approaches attempt to address this limitation. Hybrid methods combine data-driven generation with rule-based systems that filter or constrain outputs according to known medical knowledge. For example, generated electronic health records can be validated against clinical guidelines that specify valid medication combinations and contraindicated prescription patterns. Mendes et al. emphasized the importance of incorporating domain expertise throughout the generation pipeline rather than treating validation as a post-hoc step \cite{mendes2025synthetic}.

Adversarial validation provides another quality assurance mechanism, where a separate classifier attempts to distinguish synthetic from real samples. If the discriminator achieves high accuracy, this indicates that synthetic samples exhibit detectable artifacts or systematic differences from real data. However, this approach fundamentally measures statistical similarity rather than biological plausibility---a model might generate statistically indistinguishable samples that nonetheless violate medical constraints not represented in the training distribution.

The computational cost of synthetic data generation also warrants consideration in resource allocation decisions. Finetti et al. noted that deep generative model training typically requires GPU clusters and substantial hyperparameter tuning effort \cite{finetti2025data}. For datasets with fewer than 100 samples, the overhead of GAN training may exceed the benefit relative to simpler augmentation techniques. The decision between classical augmentation, VAEs, GANs, or diffusion models must weigh not only performance metrics but also practical constraints on computational resources, domain expertise availability, and time-to-deployment requirements.
